% GENERATED FILE. Edit canonical lessons, then run npm run generate:books.
% canonical-source-hash: fnv1a64:a5b5563e85db83ed
% canonical-lessons: FR-C142-lhoraire, FR-C142-lavis, FR-C142-le-loyer, FR-C142-lannonce, FR-C142-la-facture, FR-R142-first-pass-saying-no-and-looking-back, FR-R142-second-pass-looking-ahead-and-reading-notices

\chapter{Opening Hours, a Notice, the Rent, an Advert, a Bill}
\label{ch:fr-opening-hours-a-notice-the-rent-an-advert-a-bill}
\hlchaptermodality{142}

\begin{chapteropening}
\textbf{By the end of this chapter:} \emph{I can read opening hours, a notice, an advert and a bill, and act on what they say.}

Complete the last lesson of chapter 142: I can read opening hours, a notice, an advert and a bill, and act on what they say.
\end{chapteropening}

\section[l'horaire]{l'horaire — opening hours, a timetable}

\label{lesson:FR-C142-lhoraire}



\begin{quote}
Before the new one: say the French for to foresee, then the French for to organise.
\end{quote}

\subsection*{You'll want to know: l'horaire}
\textbf{l'horaire} — \textquotedblleft{}opening hours, a timetable\textquotedblright{}.

\emph{Read it:} the door

The door says \textbf{Horaire : lundi – vendredi, 9 h – 18 h.} Open on weekdays from nine to six: not on \textbf{samedi}.

\begin{grammarlens}[title={What you've built}]
One more word for this chapter.
\end{grammarlens}

\subsection*{Your turn}
\begin{itemize}
\raggedright
  \item \emph{Say it:} \emph{l'horaire}
  \item \emph{Say it:} \emph{l'horaire}, once more
  \item \emph{Say it:} say \emph{l'horaire}
  \item \emph{Recall:} say the French for to foresee, then the French for to organise, then say \emph{l'horaire} again
\end{itemize}

\begin{tcolorbox}[breakable,colback=teal!4,colframe=teal!35!black,title={Before you move on}]
What does l'horaire mean? (\textquotedblleft{}Opening hours, a timetable\textquotedblright{}.) Say it once more.
\end{tcolorbox}
\section[l'avis]{l'avis — a notice, an opinion}

\label{lesson:FR-C142-lavis}



\begin{quote}
Before the new one: say the French for to organise, then the French for opening hours.
\end{quote}

\subsection*{You'll want to know: l'avis}
\textbf{l'avis} — \textquotedblleft{}a notice, an opinion\textquotedblright{}.

\emph{Read it:} the notice

The notice says \textbf{Avis : le magasin est fermé aujourd'hui.} The shop is closed today, so come back another day.

\begin{grammarlens}[title={What you've built}]
One more word for this chapter.
\end{grammarlens}

\subsection*{Your turn}
\begin{itemize}
\raggedright
  \item \emph{Say it:} \emph{l'avis}
  \item \emph{Say it:} \emph{l'avis}, once more
  \item \emph{Say it:} say \emph{l'avis}
  \item \emph{Recall:} say the French for to organise, then the French for opening hours, then say \emph{l'avis} again
\end{itemize}

\begin{tcolorbox}[breakable,colback=teal!4,colframe=teal!35!black,title={Before you move on}]
What does l'avis mean? (\textquotedblleft{}A notice, an opinion\textquotedblright{}.) Say it once more.
\end{tcolorbox}
\section[le loyer]{le loyer — the rent}

\label{lesson:FR-C142-le-loyer}



\begin{quote}
Before the new one: say the French for opening hours, then the French for a notice.
\end{quote}

\subsection*{You'll want to know: le loyer}
\textbf{le loyer} — \textquotedblleft{}the rent\textquotedblright{}. The money for a \textbf{chambre}, paid every \textbf{mois}: \textbf{Le loyer est de 400 euros par mois.}

\begin{grammarlens}[title={What you've built}]
One more word for this chapter.
\end{grammarlens}

\subsection*{Your turn}
\begin{itemize}
\raggedright
  \item \emph{Say it:} \emph{le loyer}
  \item \emph{Say it:} \emph{le loyer}, once more
  \item \emph{Say it:} say \emph{le loyer}
  \item \emph{Recall:} say the French for opening hours, then the French for a notice, then say \emph{le loyer} again
\end{itemize}

\begin{tcolorbox}[breakable,colback=teal!4,colframe=teal!35!black,title={Before you move on}]
What does le loyer mean? (\textquotedblleft{}The rent\textquotedblright{}.) Say it once more.
\end{tcolorbox}
\section[l'annonce]{l'annonce — an advert, an announcement}

\label{lesson:FR-C142-lannonce}



\begin{quote}
Before the new one: say the French for a notice, then the French for the rent.
\end{quote}

\subsection*{You'll want to know: l'annonce}
\textbf{l'annonce} — \textquotedblleft{}an advert, an announcement\textquotedblright{}.

\emph{Read it:} the advert

The advert says \textbf{Chambre libre. Loyer : 400 euros.} A room is free, at 400 a month.

\begin{grammarlens}[title={What you've built}]
One more word for this chapter.
\end{grammarlens}

\subsection*{Your turn}
\begin{itemize}
\raggedright
  \item \emph{Say it:} \emph{l'annonce}
  \item \emph{Say it:} \emph{l'annonce}, once more
  \item \emph{Say it:} say \emph{l'annonce}
  \item \emph{Recall:} say the French for a notice, then the French for the rent, then say \emph{l'annonce} again
\end{itemize}

\begin{tcolorbox}[breakable,colback=teal!4,colframe=teal!35!black,title={Before you move on}]
What does l'annonce mean? (\textquotedblleft{}An advert, an announcement\textquotedblright{}.) Say it once more.
\end{tcolorbox}
\section[la facture]{la facture — a bill, an invoice}

\label{lesson:FR-C142-la-facture}



\begin{quote}
Before the new one: say the French for the rent, then the French for an advert.
\end{quote}

\subsection*{You'll want to know: la facture}
\textbf{la facture} — \textquotedblleft{}a bill, an invoice\textquotedblright{}.

\emph{Read it:} the bill

The bill says \textbf{Facture : loyer, 400 euros. À payer avant lundi.} Pay it before Monday.

\begin{grammarlens}[title={What you've built}]
One more word for this chapter.
\end{grammarlens}

\subsection*{Your turn}
\begin{itemize}
\raggedright
  \item \emph{Say it:} \emph{la facture}
  \item \emph{Say it:} \emph{la facture}, once more
  \item \emph{Say it:} say \emph{la facture}
  \item \emph{Recall:} say the French for the rent, then the French for an advert, then say \emph{la facture} again
\end{itemize}

\begin{tcolorbox}[breakable,colback=teal!4,colframe=teal!35!black,title={Before you move on}]
What does la facture mean? (\textquotedblleft{}A bill, an invoice\textquotedblright{}.) Say it once more.
\end{tcolorbox}
\section[(dialogue)]{Saying no and looking back — a first pass}

\label{lesson:FR-R142-first-pass-saying-no-and-looking-back}



\begin{quote}
Say the words for an advert and a bill.
\end{quote}

\subsection*{Your turn}
Say these aloud:

\begin{itemize}
\raggedright
  \item nothing, none, neither, which one, anything at all
  \item already, recently, the day before yesterday, formerly, to happen
\end{itemize}

\begin{tcolorbox}[breakable,colback=teal!4,colframe=teal!35!black,title={Before you move on}]
Nothing? (\textbf{rien}.) None? (\textbf{aucun / aucune}.) Neither? (\textbf{non plus}.) Which one? (\textbf{lequel / laquelle}.) Anything at all? (\textbf{n'importe quoi}.) Already? (\textbf{déjà}.) Recently? (\textbf{récemment}.) The day before yesterday? (\textbf{avant-hier}.) Formerly? (\textbf{autrefois}.) To happen? (\textbf{se produire}.)
\end{tcolorbox}
\section[(dialogue)]{Looking ahead and reading notices — a second pass}

\label{lesson:FR-R142-second-pass-looking-ahead-and-reading-notices}



\begin{quote}
Say the words for the future and this evening.
\end{quote}

\subsection*{Your turn}
Say these aloud:

\begin{itemize}
\raggedright
  \item the future, this evening, a plan, to foresee, to organise
  \item opening hours, a notice, the rent, an advert, a bill
\end{itemize}

\begin{tcolorbox}[breakable,colback=teal!4,colframe=teal!35!black,title={Before you move on}]
The future? (\textbf{l'avenir}.) This evening? (\textbf{ce soir}.) A plan? (\textbf{le projet}.) To foresee? (\textbf{prévoir}.) To organise? (\textbf{organiser}.) Opening hours? (\textbf{l'horaire}.) A notice? (\textbf{l'avis}.) The rent? (\textbf{le loyer}.) An advert? (\textbf{l'annonce}.) A bill? (\textbf{la facture}.)
\end{tcolorbox}
